\section{Python-Implementierung von Markov-Ketten}

Markov-Ketten lassen sich mit Python sehr elegant implementieren und visualisieren. Im Folgenden zeigen wir eine einfache Implementation, die es erlaubt, Markov-Ketten zu simulieren und deren Eigenschaften zu analysieren.

\begin{lstlisting}[language=Python, caption=Grundlegende Implementation einer Markov-Kette in Python]
import numpy as np
import matplotlib.pyplot as plt

# Einfache Markov-Ketten Implementierung
class MarkovChain:
    def __init__(self, transition_matrix, states=None):
        """
        Initialisiert eine Markov-Kette mit einer Übergangsmatrix.

        Args:
            transition_matrix: Eine quadratische Matrix als numpy-Array
            states: Optional - Liste mit Namen der Zustände
        """
        self.P = np.array(transition_matrix)
        n = self.P.shape[0]
        if states is None:
            self.states = [str(i) for i in range(n)]
        else:
            self.states = states

        # Überprüfen, ob die Übergangsmatrix gültig ist
        row_sums = self.P.sum(axis=1)
        if not np.allclose(row_sums, 1.0):
            raise ValueError("Die Zeilen der Übergangsmatrix müssen sich zu 1 summieren.")

    def simulate(self, steps, start_state=0):
        """
        Simuliert die Markov-Kette für eine bestimmte Anzahl von Schritten.

        Args:
            steps: Anzahl der Simulationsschritte
            start_state: Anfangszustand (als Index)

        Returns:
            Liste der besuchten Zustände (als Indizes)
        """
        n = self.P.shape[0]
        if start_state >= n:
            raise ValueError(f"Startzustand muss kleiner als {n} sein")

        # Simulation initialisieren
        states = [start_state]
        current = start_state

        # Für jeden Schritt
        for _ in range(steps):
            # Übergangswahrscheinlichkeiten für den aktuellen Zustand
            probs = self.P[current]

            # Nächsten Zustand zufällig auswählen
            current = np.random.choice(n, p=probs)
            states.append(current)

        return states

    def compute_stationary_distribution(self):
        """
        Berechnet die stationäre Verteilung der Markov-Kette.

        Returns:
            Stationäre Verteilung als numpy-Array
        """
        # Eigenwerte und Eigenvektoren der transponierten Matrix
        eigenvalues, eigenvectors = np.linalg.eig(self.P.T)

        # Finde den Eigenvektor zum Eigenwert 1
        one_eigenvector = eigenvectors[:, np.isclose(eigenvalues, 1.0)]

        # Normalisieren, damit die Summe 1 ergibt
        stationary = one_eigenvector[:, 0].real
        stationary = stationary / stationary.sum()

        return stationary
\end{lstlisting}

Mit dieser Klasse können wir nun einfach Markov-Ketten simulieren und analysieren. Zum Beispiel können wir unser Wetter-Beispiel implementieren:

\begin{lstlisting}[language=Python, caption=Anwendung der MarkovChain-Klasse auf das Wetter-Beispiel]
# Wetter-Markov-Kette
weather_matrix = [
    [0.8, 0.2],  # Sonnig -> Sonnig, Sonnig -> Regnerisch
    [0.4, 0.6]   # Regnerisch -> Sonnig, Regnerisch -> Regnerisch
]

weather_chain = MarkovChain(weather_matrix, states=["Sonnig", "Regnerisch"])

# Simulation für 20 Tage, beginnend mit sonnigem Wetter
simulation = weather_chain.simulate(20, start_state=0)
weather_states = [weather_chain.states[s] for s in simulation]

print("Wettersimulation für 20 Tage:")
print(weather_states)

# Stationäre Verteilung berechnen
stationary = weather_chain.compute_stationary_distribution()
print("\nStationäre Verteilung:")
for i, state in enumerate(weather_chain.states):
    print(f"{state}: {stationary[i]:.2f}")
\end{lstlisting}

Eine mögliche Ausgabe könnte sein:
\begin{verbatim}
Wettersimulation für 20 Tage:
['Sonnig', 'Sonnig', 'Sonnig', 'Regnerisch', 'Regnerisch', 'Sonnig', 
'Sonnig', 'Sonnig', 'Sonnig', 'Sonnig', 'Sonnig', 'Regnerisch', 
'Regnerisch', 'Regnerisch', 'Sonnig', 'Sonnig', 'Sonnig', 'Sonnig', 
'Sonnig', 'Sonnig', 'Sonnig']

Stationäre Verteilung:
Sonnig: 0.67
Regnerisch: 0.33
\end{verbatim}

Diese Python-Implementation erlaubt es uns, schnell verschiedene Markov-Ketten zu modellieren und zu simulieren. Besonders nützlich ist die Möglichkeit, die stationäre Verteilung zu berechnen und zu visualisieren, wie sich die Zustandsverteilung über die Zeit entwickelt.

\begin{myexercise}[Python-Simulation]
    Erweitern Sie den obigen Code, um die App-Nutzung aus unserem früheren Beispiel zu simulieren. Die Übergangsmatrix ist:
    \[
        P = \begin{pmatrix}
            0{,}7 & 0{,}3 \\
            0{,}3 & 0{,}7 \\
        \end{pmatrix}
    \]

    Führen Sie eine Simulation für 100 Schritte durch und berechnen Sie die stationäre Verteilung. Vergleichen Sie die stationäre Verteilung mit dem analytischen Ergebnis aus dem früheren Beispiel.

    \begin{myanswer}
        Hier ist der Code für die Simulation der App-Nutzung:

        \begin{lstlisting}[language=Python]
# App-Nutzungs-Markov-Kette
app_matrix = [
    [0.7, 0.3],  # WhatsApp -> WhatsApp, WhatsApp -> Instagram
    [0.3, 0.7]   # Instagram -> WhatsApp, Instagram -> Instagram
]

app_chain = MarkovChain(app_matrix, states=["WhatsApp", "Instagram"])

# Simulation für 100 Nutzungen, beginnend mit WhatsApp
simulation = app_chain.simulate(100, start_state=0)

# Zählen, wie oft jede App genutzt wurde
whatsapp_count = simulation.count(0)
instagram_count = simulation.count(1)
print(f"WhatsApp wurde {whatsapp_count} mal genutzt")
print(f"Instagram wurde {instagram_count} mal genutzt")
print(f"Verhältnis WhatsApp/Instagram: {whatsapp_count/(whatsapp_count+instagram_count):.2f}")

# Stationäre Verteilung berechnen
stationary = app_chain.compute_stationary_distribution()
print("\nStationäre Verteilung:")
for i, state in enumerate(app_chain.states):
    print(f"{state}: {stationary[i]:.2f}")
\end{lstlisting}

        Die stationäre Verteilung ist $\pi = (0.5, 0.5)$, was genau dem analytischen Ergebnis entspricht. Bei einer langen Simulation sollte sich das Verhältnis der App-Nutzung diesem Wert annähern.
    \end{myanswer}
\end{myexercise}

\section{PageRank-Algorithmus: Eine Anwendung von Markov-Ketten}

Eine der bekanntesten Anwendungen von Markov-Ketten ist der PageRank-Algorithmus, der von Google-Gründern Larry Page und Sergey Brin entwickelt wurde. PageRank ist ein Algorithmus zur Bewertung der Wichtigkeit von Webseiten und bildete die Grundlage der frühen Google-Suchmaschine.

\subsection{Grundprinzip des PageRank}

Die grundlegende Idee hinter PageRank ist einfach und elegant: Die Wichtigkeit einer Webseite wird durch die Anzahl und Qualität der Links bestimmt, die auf sie verweisen. Eine Webseite ist wichtig, wenn viele andere wichtige Seiten auf sie verlinken.

Diese rekursive Definition lässt sich mathematisch als Markov-Kette modellieren:

\begin{itemize}
    \item Jede Webseite entspricht einem Zustand in der Markov-Kette.
    \item Die Übergangswahrscheinlichkeiten werden aus der Linkstruktur abgeleitet: Wenn Seite A auf $n$ andere Seiten verlinkt, beträgt die Übergangswahrscheinlichkeit zu jeder dieser Seiten $1/n$.
    \item Der PageRank einer Seite ist proportional zur Wahrscheinlichkeit, dass ein zufällig durch das Web surfender Nutzer auf dieser Seite landet.
\end{itemize}

\subsection{Mathematische Formulierung}

Formal lässt sich der PageRank-Wert $PR(p_i)$ einer Seite $p_i$ wie folgt definieren:

\[
    PR(p_i) = \frac{1-d}{N} + d \sum_{p_j \in M(p_i)} \frac{PR(p_j)}{L(p_j)}
\]

wobei:
\begin{itemize}
    \item $d$ ist ein Dämpfungsfaktor, typischerweise $0{,}85$
    \item $N$ ist die Gesamtzahl der Webseiten
    \item $M(p_i)$ ist die Menge aller Seiten, die auf $p_i$ verlinken
    \item $L(p_j)$ ist die Anzahl der ausgehenden Links von Seite $p_j$
\end{itemize}

Der Dämpfungsfaktor $d$ modelliert die Wahrscheinlichkeit, dass ein Nutzer einem Link folgt, anstatt direkt zu einer zufälligen anderen Seite zu springen.

\subsection{Beziehung zu Markov-Ketten}

Der PageRank-Algorithmus lässt sich direkt als Markov-Kette interpretieren:

\begin{itemize}
    \item Die Übergangsmatrix $P$ beschreibt die Linkstruktur des Webs.
    \item Der PageRank-Vektor entspricht der stationären Verteilung dieser Markov-Kette.
    \item Der Dämpfungsfaktor sorgt dafür, dass die Markov-Kette irreduzibel und aperiodisch wird, was die Existenz und Eindeutigkeit der stationären Verteilung garantiert.
\end{itemize}

\begin{myexample}[PageRank für ein kleines Web]
Betrachten wir ein einfaches Web mit vier Seiten A, B, C und D mit folgender Linkstruktur:
\begin{itemize}
    \item Seite A verlinkt auf B und C
    \item Seite B verlinkt auf C
    \item Seite C verlinkt auf A
    \item Seite D verlinkt auf A, B und C
\end{itemize}

\begin{figure}[H]
    \centering
    \begin{tikzpicture}[node distance=2.5cm, auto, thick]
        \node[state] (A) {A};
        \node[state] (B) [right of=A] {B};
        \node[state] (C) [below of=B] {C};
        \node[state] (D) [below of=A] {D};

        \path[->]
        (A) edge (B)
        (A) edge (C)
        (B) edge (C)
        (C) edge (A)
        (D) edge (A)
        (D) edge (B)
        (D) edge (C);
    \end{tikzpicture}
    \caption{Linkstruktur des Beispiel-Webs}
    \label{fig:pagerank-example}
\end{figure}

Die Übergangsmatrix ohne Dämpfungsfaktor wäre:
\[
    P = \begin{pmatrix}
        0    & 0.5  & 0.5  & 0 \\
        0    & 0    & 1    & 0 \\
        1    & 0    & 0    & 0 \\
        0.33 & 0.33 & 0.33 & 0
    \end{pmatrix}
\]

Mit einem Dämpfungsfaktor von $d = 0{,}85$ und $N = 4$ wird die angepasste Übergangsmatrix:
\[
    P' = 0{,}85 \cdot P + 0{,}15 \cdot \frac{1}{4} \cdot \begin{pmatrix}
        1 & 1 & 1 & 1 \\
        1 & 1 & 1 & 1 \\
        1 & 1 & 1 & 1 \\
        1 & 1 & 1 & 1
    \end{pmatrix}
\]

Die stationäre Verteilung dieser Markov-Kette, berechnet durch Lösen von $\pi = \pi \cdot P'$ und $\sum \pi_i = 1$, ergibt den PageRank-Vektor:
\[
    PR \approx (0.32, 0.17, 0.39, 0.12)
\]

Dies bedeutet:
\begin{itemize}
    \item Seite C hat den höchsten PageRank (0.39)
    \item Gefolgt von Seite A (0.32)
    \item Dann Seite B (0.17)
    \item Seite D hat den niedrigsten PageRank (0.12)
\end{itemize}

Dieses Ergebnis ist intuitiv nachvollziehbar: Seite C erhält Links von drei Seiten (A, B und D), während Seite D keine eingehenden Links hat.
\end{myexample}

\subsection{PageRank-Implementierung in Python}

Mit unserer MarkovChain-Klasse können wir eine einfache Version des PageRank-Algorithmus implementieren:

\begin{lstlisting}[language=Python, caption=Einfache PageRank-Implementierung mit der MarkovChain-Klasse]
def calculate_pagerank(links, damping_factor=0.85, num_pages=None):
    """
    Berechnet PageRank für ein Netzwerk von Webseiten.

    Args:
        links: Dictionary mit {Seite: [Liste der Seiten, auf die verlinkt wird]}
        damping_factor: Dämpfungsfaktor (Standard: 0.85)
        num_pages: Anzahl der Seiten (Optional, wird sonst automatisch bestimmt)

    Returns:
        Dictionary mit {Seite: PageRank-Wert}
    """
    # Alle Seiten identifizieren (verlinkte und verlinkende)
    all_pages = set()
    for page, outlinks in links.items():
        all_pages.add(page)
        all_pages.update(outlinks)

    # Seiten in sortierte Liste umwandeln für konsistente Indizierung
    pages = sorted(list(all_pages))
    n = len(pages)

    if num_pages is not None and num_pages > n:
        n = num_pages

    # Übergangsmatrix erstellen
    P = np.zeros((n, n))

    # Übergangswahrscheinlichkeiten basierend auf Links setzen
    page_indices = {page: i for i, page in enumerate(pages)}
    for page, outlinks in links.items():
        if outlinks:  # Wenn die Seite ausgehende Links hat
            i = page_indices[page]
            for target in outlinks:
                j = page_indices[target]
                P[i, j] = 1.0 / len(outlinks)

    # Dämpfungsfaktor anwenden
    P = damping_factor * P + (1 - damping_factor) / n * np.ones((n, n))

    # Markov-Kette erstellen und stationäre Verteilung berechnen
    mc = MarkovChain(P, states=pages)
    stationary = mc.compute_stationary_distribution()

    # PageRank als Dictionary zurückgeben
    pagerank = {page: stationary[i] for i, page in enumerate(pages)}
    return pagerank

# Beispiel: Unser kleines Web
web = {
    'A': ['B', 'C'],
    'B': ['C'],
    'C': ['A'],
    'D': ['A', 'B', 'C']
}

pagerank = calculate_pagerank(web)
for page, rank in sorted(pagerank.items()):
    print(f"Seite {page}: PageRank = {rank:.4f}")
\end{lstlisting}

Der PageRank-Algorithmus zeigt eindrücklich, wie Markov-Ketten in der Praxis eingesetzt werden können, um komplexe Probleme zu lösen. Die stationäre Verteilung einer geeignet konstruierten Markov-Kette kann wertvolle Einblicke in die Struktur vernetzter Systeme liefern.